\chapter{Conclusion and Outlook}
Theoretical chemistry is defined by the constant tension between accuracy and computational cost. In this thesis, I have presented an overview of my efforts to develop and apply methods that strike a balance between these two competing demands. A significant portion of my work has been dedicated to the development of the \textsc{Zinq} software package, that allows us to easily integrate, validate, and apply a wide range of quantum dynamical methods. This software has been instrumental in my research, enabling me to explore and compare various approaches to quantum dynamics.

With this robust toolkit, I have pushed the boundaries of surface hopping methods to evaluate their limits and identify areas for improvement. My Exploration of curvature-driven surface hopping algorithms demonstrated that such shortcuts have their boundaries. Bypassing the need for \glspl{nacv}, these methods can be computationally efficient, but they may not always capture the full complexity of nonadiabatic processes. My work has highlighted the importance of carefully assessing the applicability of these methods by comparing them against more rigorous approaches, such as \gls{fssh}. Through systematic benchmarking and analysis, I have identified specific scenarios where curvature-driven algorithms may fall short, providing valuable insights for future method development. This work has not only advanced our understanding of the strengths and limitations of approximating \gls{nacv} but also paved the way for the development of more accurate and efficient algorithms in the future.

While algorithmic rigor in quantum dynamics is crucial, my sensitivity analysis of \textit{cis}-stilbene photodynamics has revealed that the choice of electronic structure method is equally important. I showed that relying solely on short-time observables, such as \gls{trpes}, is insufficient to correctly validate the method's ability to predict long-time photoproduct quantum yields. Because even minor energetic shifts can lead to vastly different photoproduct distributions, this this work serves as a cautionary tale for the field, emphasizing the need for comprehensive validation of electronic structure methods in the context of nonadiabatic dynamics.

Beyond molecular photochemistry, this fork successfully bridges the gap between theoretical chemistry and collision physics. By applying quantum dynamics methods to the study of ion-atom collisions, I have demonstrated the versatility and broad applicability of these techniques. This work has not only provided insights into the spin-orbit induced charge transfer processes governing ion-atom interactions but also highlighted the potential for cross-disciplinary research that can benefit both fields.

Looking ahead, there are several exciting avenues for future research. I plan to continue refining and expanding the capabilities of the \textsc{Zinq} software, incorporating new algorithms and features that will further enhance its utility for the research community. Additionally, the ability of surface hopping algorithms to accurately capture nonadiabatic dynamics remains an active area of investigation. I intend to explore new approaches to improve the accuracy of these methods, particularly in challenging scenarios where traditional approximations may fail. Finally, I am eager to further explore the intersection of theoretical chemistry and collision physics, applying advanced quantum dynamics methods to a wider range of systems and phenomena, with the goal of improving the accuracy and predictive power of our models in both fields.
